\documentclass[10pt,answers]{exam}
\usepackage[legalpaper,bindingoffset=0in,left=.6in,right=0.8in,top=0.4in,bottom=0.7in,headsep=.5\baselineskip]{geometry}
\usepackage[utf8]{inputenc}
\usepackage{graphicx}
\usepackage{listings}
\usepackage{enumitem}
\usepackage{amsmath}
\usepackage{multirow}
\usepackage{multicol}
\usepackage{caption}
\usepackage{tabularx}
\usepackage{url}
\usepackage{amssymb}
\usepackage{array}
\usepackage{tikz}
\usetikzlibrary{shapes.callouts, positioning, arrows.meta, shapes.geometric, shadows, calc, patterns, 3d, backgrounds, shadings, decorations.markings}
\usepackage{adjustbox}
\usepackage{ragged2e}
\usepackage{pgfplots}
\pgfplotsset{compat=1.18}

\def\changemargin#1#2{\list{}{\rightmargin#2\leftmargin#1}
\item[]}
\let\endchangemargin=\endlist

\usepackage{xcolor}
\definecolor{mGreen}{rgb}{0,0.6,0}
\definecolor{mGray}{rgb}{0.5,0.5,0.5}
\definecolor{mPurple}{rgb}{0.58,0,0.82}
\definecolor{backgroundColour}{rgb}{0.95,0.95,0.92}
\lstset{
  language=C,
  basicstyle=\ttfamily\small,
  aboveskip={1.0\baselineskip},
  belowskip={1.0\baselineskip},
  columns=fixed,
  extendedchars=true,
  breaklines=true,
  tabsize=4,
  prebreak=\raisebox{0ex}[0ex][0ex]{\ensuremath{\hookleftarrow}},
  frame=lines,
  showtabs=false,
  showspaces=false,
  showstringspaces=false,
  keywordstyle=\color[rgb]{0.627,0.126,0.941},
  commentstyle=\color[rgb]{0.133,0.545,0.133},
  stringstyle=\color[rgb]{01,0,0},
  numbers=left,
  numberstyle=\small,
  stepnumber=1,
  numbersep=10pt,
  captionpos=t,
escapeinside={\%*}{*)}
}

\pointsdroppedatright

\title{
\vspace{-1.5cm}
\begin{changemargin}{0.68cm}{0.68cm}
  \normalsize Class Test 4 Makeup Solutions\hfill {July 2025}\\
\end{changemargin}
\vspace{.4cm}
\large BANGLADESH UNIVERSITY OF ENGINEERING AND TECHNOLOGY\\
\vspace{.1cm}
Department of Computer Science and Engineering\\
\vspace{.2cm}
Course Code: CSE 219 \hspace{.1cm} Course Title: Signals and Linear Systems\\
\vspace{.2cm}
\textbf{Time: 25 minutes} \hspace{1.7cm} \textbf{Marks: 16 + 4 bonus}
\begin{changemargin}{0.68cm}{0.68cm}
  \normalsize Student Name: \underline{Md. Ashrafur Rahman Khan} \hfill Student ID: \underline{1905005}\\
\end{changemargin}
}
\date{}

\begin{document}

\maketitle
\thispagestyle{empty}
\vspace{-2.4cm}

\begin{center}
Answer \textbf{all of Questions 1--2}. \textbf{Question 3 is for bonus marks.} Answer concisely.
\end{center}

\begin{questions}

% Question 1
\question Consider the sequence
\[
  x[n] = \{2,0,1,-1\}, \qquad 0 \le n < 4
\]
You want to compute its 4-point DFT using the \textbf{radix-2 DIF FFT}. \hfill [8]

\begin{parts}
  \part Perform the first butterfly stage and write the intermediate values.
  \begin{solution}
    In radix-2 DIF FFT, the first stage combines the pairs $(x[0],x[2])$ and $(x[1],x[3])$.

    From the first pair:
    \[
      a_0 = x[0] + x[2] = 2 + 1 = 3,
    \]
    \[
      a_2 = (x[0] - x[2])W_4^0 = (2-1)\cdot 1 = 1.
    \]

    From the second pair:
    \[
      a_1 = x[1] + x[3] = 0 + (-1) = -1,
    \]
    \[
      a_3 = (x[1] - x[3])W_4^1 = (0-(-1))W_4^1 = 1\cdot(-j) = -j.
    \]

    So after the first stage, the intermediate values are
    \[
      \{a_0,a_1,a_2,a_3\} = \{3,-1,1,-j\}.
    \]
  \end{solution}
  \part Complete the second stage and write the outputs in the order produced directly by the DIF algorithm.
  \begin{solution}
    The second stage applies 2-point butterflies separately to $(a_0,a_1)$ and $(a_2,a_3)$:
    \[
      y_0 = a_0 + a_1 = 3 + (-1) = 2,
    \]
    \[
      y_1 = a_0 - a_1 = 3 - (-1) = 4,
    \]
    \[
      y_2 = a_2 + a_3 = 1 + (-j) = 1-j,
    \]
    \[
      y_3 = a_2 - a_3 = 1 - (-j) = 1+j.
    \]
    Therefore, the outputs produced directly by the DIF flow are
    \[
      \{y_0,y_1,y_2,y_3\} = \{2,4,1-j,1+j\}.
    \]
  \end{solution}
  \part Rearrange the result from part (b) to obtain the DFT in natural order, i.e. $X[0], X[1], X[2], X[3]$.
  \begin{solution}
    For radix-2 DIF FFT, the outputs come out in bit-reversed order. For $N=4$, the bit-reversal mapping is
    \[
      0 \to 0,\qquad 1 \to 2,\qquad 2 \to 1,\qquad 3 \to 3.
    \]
    So the algorithm-order output
    \[
      \{2,4,1-j,1+j\}
    \]
    corresponds to
    \[
      \{X[0],X[2],X[1],X[3]\}.
    \]
    Hence the DFT in natural order is
    \[
      X[0]=2,\qquad X[1]=1-j,\qquad X[2]=4,\qquad X[3]=1+j.
    \]
  \end{solution}
  \part Briefly state where bit-reversal appears in radix-2 DIT FFT and where it appears in radix-2 DIF FFT.
  \begin{solution}
    In radix-2 DIT FFT, bit-reversal appears at the \textbf{input} side: the algorithm usually expects the inputs in bit-reversed order and produces outputs in natural order.

    In radix-2 DIF FFT, bit-reversal appears at the \textbf{output} side: the inputs are in natural order, and the algorithm produces outputs in bit-reversed order.
  \end{solution}
\end{parts}
\vspace{0.2cm}

% Question 2
\question In an audio-processing application, a discrete-time signal $x[n]$ of length $L$ is passed through an FIR filter $h[n]$ of length $M$ to produce
\[
  y[n] = (x*h)[n].
\]
For large $L$ and $M$, you want to compute this efficiently. \hfill [8]

\begin{parts}
  \part What is the length of the output sequence $y[n]$? What is the time complexity of direct time-domain computation?
  \begin{solution}
    If $x[n]$ has length $L$ and $h[n]$ has length $M$, then their linear convolution has length
    \[
      L+M-1.
    \]

    A direct implementation computes, for each output sample, a sum over filter taps. This leads to time complexity
    \[
      O(LM).
    \]
  \end{solution}
  \part If you compute DFTs of insufficient length, multiply them, and then take an IDFT, what kind of convolution do you obtain instead of linear convolution?
  \begin{solution}
    You obtain \textbf{circular convolution}, not linear convolution.

    In the audio context, this causes wrap-around: samples from the tail of the filtered signal fold back and corrupt the beginning of the output. So the result is not the correct FIR-filtered signal.
  \end{solution}
  \part Propose an FFT-based algorithm to compute the correct filtered output.
  \begin{solution}
    A correct FFT-based method is:
            \begin{lstlisting}
function FastFiltering(x, h):
    Choose N
    Zero-pad x and h to length N
    X = FFT(x_padded)
    H = FFT(h_padded)
    Y = X .* H              // pointwise multiplication
    y_padded = IFFT(Y)
    Return the first L+M-1 samples
\end{lstlisting}
    This works because multiplication in the frequency domain corresponds to convolution in the time domain, and the zero-padding prevents wrap-around from corrupting the result.
  \end{solution}
  \part If the transform length is $N$, what condition must $N$ satisfy? What is the resulting asymptotic complexity?
  \begin{solution}
    To make circular convolution equal linear convolution, the transform length must satisfy
    \[
      N \ge L + M - 1.
    \]
    In practice, we often choose the next convenient FFT length, often a power of 2.

    The algorithm uses two FFTs, one pointwise multiplication, and one inverse FFT. Therefore its asymptotic complexity is
    \[
      O(N \log N).
    \]
    This is much better than the direct $O(LM)$ method for large signals.
  \end{solution}
\end{parts}
\vspace{0.2cm}

% Question 3
\question[4] Thank you for attending the classes and participating in this CT. \hfill [4]
\begin{solution}
  You're welcome!
\end{solution}

\end{questions}

\end{document}
